\section{What does it mean to understand a language model?}
\label{sec:intro}

\subsection{The gap between knowing the weights and knowing the algorithm}

A large language model (LLM) takes a piece of text and returns a probability distribution over the next token.
Inside, it is a long but completely explicit computation: matrix products, additions, a few simple nonlinear
functions. Nothing is hidden from us. We can print any weight and record any intermediate number.

And yet we cannot say how the model works in the sense in which we can say how a sorting algorithm works. The
weights were not written by a programmer. They were found by gradient descent on a prediction loss, and the
algorithm they implement is something we have to discover after the fact. The situation resembles holding the
compiled binary of a program without its source code. Every instruction is visible, and the logic is not.

\emph{Interpretability} is the study of how to recover that logic. We use the following working definition.

\begin{definition}[Interpretability]\label{def:interp}
An \emph{interpretation} of a model is a description of its internal computation that is (i) understandable by
a human and (ii) \emph{faithful}, which means that the description correctly predicts what the model computes,
including what it computes when we intervene on its internals.
\emph{Mechanistic interpretability} is the part of the field that aims at descriptions on the level of
algorithms: which quantities the model represents, where, and through which operations they are combined.
\end{definition}

The two requirements pull in opposite directions. The full list of weights is perfectly faithful and not
understandable. A one-sentence story such as ``the model looks for the subject of the sentence'' is understandable
and may be false. The whole field lives in the tension between the two, and most of its methodological care goes
into checking condition (ii).

\subsection{Why not simply ask the model?}

A modern LLM will explain its own answers when asked. These explanations are text generated by the same model,
and nothing forces them to describe the computation that produced the answer. \citet{turpin2023language} gave a
clean demonstration. They biased multiple-choice prompts, for instance by arranging the few-shot examples so that
the correct answer was always option (A). The models then often chose (A) on the test question and wrote a
step-by-step justification of it that never mentioned the pattern that had caused the choice.

The literature separates two properties of an explanation \citep{jacovi2020towards}. An explanation is
\emph{plausible} if it convinces a human reader, and it is \emph{faithful} if it reflects the true reasoning process
of the model. Self-explanations are selected for plausibility. Interpretability research is the attempt to obtain
faithfulness, and for that we have to look inside.

\subsection{Why it matters}

We give four reasons, in decreasing order of how often they are cited.

\begin{itemize}
\item \textbf{Safety and auditing.} We would like to check properties of a model that its behaviour on a test set
cannot reveal, for example whether it pursues an objective that it does not state. \citet{marks2025auditing}
trained a model with a hidden objective and let blinded teams try to find it. Tools that read internal
activations were among the ones that helped.
\item \textbf{Debugging and control.} If we know where and how a behaviour is computed, we can change it in a
targeted way. \Cref{sec:interventions} shows how adding one vector to the activations can switch a behaviour on
or off.
\item \textbf{Science.} Neural networks generalise in ways we do not fully understand. Reading off the algorithm
that a trained network implements turns questions about generalisation into questions about concrete mechanisms.
\Cref{sec:modadd} contains an example in which a small network was reverse engineered completely.
\item \textbf{Trust.} High-stakes uses of a model call for reasons, and the reasons should be the real ones.
\end{itemize}

\subsection{Three questions and three kinds of evidence}

\begin{table}[t]
\centering\small
\begin{tabular}{@{}>{\raggedright\arraybackslash}p{0.29\textwidth}>{\raggedright\arraybackslash}p{0.37\textwidth}p{0.14\textwidth}l@{}}
\toprule
Method & What it tells us & Evidence & Where \\
\midrule
Attention patterns & from which positions a head reads & observational & \cref{sec:attention} \\
Input attribution & which input tokens the output is sensitive to & observational & \cref{sec:attribution} \\
Logit lens, direct logit attribution & when, and by which component, the answer is written & observational & \cref{sec:lens} \\
Probing & what is decodable from an activation & observational & \cref{sec:probing} \\
Sparse autoencoders & a candidate dictionary of features & observational & \cref{sec:sae} \\
Ablation, activation patching & which components the behaviour depends on & interventional & \cref{sec:interventions} \\
Steering & whether a direction controls a behaviour & interventional & \cref{sec:steering} \\
Circuit analysis & the algorithm & mechanistic & \cref{sec:circuits} \\
\bottomrule
\end{tabular}
\caption{The methods covered in this tutorial. ``Observational'' methods only record activations, and
``interventional'' methods change them. A dictionary found by a sparse autoencoder is observational evidence
until its features have been tested by interventions.}
\label{tab:map}
\end{table}

Almost every paper in the field asks one of three questions about a model and a behaviour of interest.
\begin{enumerate}[label=\textbf{Q\arabic*.}, leftmargin=3em]
\item \emph{Representation.} What information is present in the activations, and in what form?
\item \emph{Localisation.} Which parts of the model matter for the behaviour?
\item \emph{Algorithm.} What computation do these parts carry out, and why does it produce the behaviour?
\end{enumerate}
The methods differ in the kind of evidence they give, and it pays to keep the kinds apart from the first day.
\emph{Observational} methods record activations on some inputs and look for structure. They are cheap, and they
establish correlations only. \emph{Interventional} methods change an activation and measure the effect on the
output. They establish causal claims, relative to the specific intervention that was made. \emph{Mechanistic}
explanations go one step further and derive the behaviour from the weights, so that the explanation also covers
inputs that were never tried. \Cref{tab:map} places the methods of this tutorial on this map.


\subsection{The running example}
\label{sec:running}

Abstract descriptions of methods are hard to digest. We therefore apply every method to the same small task, so
that the reader can see how the different kinds of evidence fit together. The task is \emph{indirect object
identification} (IOI), introduced by \citet{wang2023interpretability}. Consider the sentence
\begin{center}
\tok{When Mary and John went to the store, John gave a drink to}
\end{center}
A fluent reader continues with \tok{Mary}. To do so, one has to notice that two names were introduced, that
\tok{John} appears a second time as the subject of the second clause, and that the remaining name is the
recipient. The name that should be predicted is called the indirect object (IO), and the repeated name is called
the subject (S).

Our model is GPT-2 small \citep{radford2019language}. It is public, it runs on a laptop, and it solves this task.
We use 12 prompts that follow the template above with six pairs of names in both orders. On all 12 prompts the
most likely next token is the correct name, with an average probability of $0.48$, against $0.014$ for the
repeated name.

\begin{caution}
All numbers and data figures about GPT-2 small in this tutorial were produced by the script
\texttt{code/gpt2\_examples.py} that accompanies the text. They are meant as illustrations on 12 prompts
from one template. They agree qualitatively with \citet{wang2023interpretability}, who studied the task on a much
larger and more varied set of prompts, and they should not be read as new findings.
\end{caution}

\subsection{How to read this tutorial}

\begin{figure}[t]
\centering
\begin{tikzpicture}[>=Latex, font=\small, node distance=6mm,
  stage/.style={draw=tutgrey, rounded corners=2pt, align=center, minimum height=12mm, inner sep=5pt, fill=#1},
  stage/.default=tutpaper!70]
\node[stage] (m) {the model\\ \cref{sec:transformer}};
\node[stage=tutblue!10, right=of m] (o) {observe\\ \cref{sec:observation}};
\node[stage=tutblue!10, right=of o] (f) {decompose\\ \cref{sec:features,sec:sae}};
\node[stage=tutred!10, right=of f] (i) {intervene\\ \cref{sec:interventions}};
\node[stage=tutred!10, right=of i] (c) {explain\\ \cref{sec:circuits}};
\node[stage, right=of c] (v) {doubt\\ \cref{sec:pitfalls}};
\foreach \a/\b in {m/o,o/f,f/i,i/c,c/v} \draw[->, tutgrey] (\a) -- (\b);
\end{tikzpicture}
\caption{Road map. Blue stages produce hypotheses from recorded activations, and red stages test them by changing
activations. The last stage collects the ways in which each of the earlier ones can mislead.}
\label{fig:roadmap}
\end{figure}

The text assumes linear algebra (vectors, matrices, inner products, projections), basic probability and the idea
of training a neural network by gradient descent. It does not assume any knowledge of transformers, which
\cref{sec:transformer} introduces from scratch. The order of the sections follows \cref{fig:roadmap}.

Three kinds of boxes mark material with a special role. An \emph{Intuition} box gives an informal picture of a
formal statement. A \emph{Caution} box states a limit of a method or a common misreading. A \emph{Takeaways} box
closes each section. Symbols are defined at the place where they are first used. Proofs are written as chains of
equalities in which the reason for each line is printed in grey on the right. Each section ends with a few
exercises, and \cref{app:solutions} contains solutions. A glossary of the technical terms is in
\cref{app:glossary}. Finally, \cref{sec:practice} lists software, first projects and further reading.
